\chapter{Introduction}

Unification is the process of trying to make two expressions of some language equal by instantiating their variables appropriately. Depending on the language of the expressions there might be multiple solutions and a special solution, the most general unifier, which "includes" all other possible solutions (see Section \ref{sec:unification} for more detail). On the other hand, there might not be a solution at all if the two terms are not unifiable.

Unification is used in a variety of applications. One example is type inference where unification can be used to find a solution to a set of constraints that were gathered to determine the most general type of some expression. Another example are proof assistants where unification can be used to determine if a rule can be applied to a set of terms. In a similar way, unification can be used to find out if a rule of a rewrite system matches a term.

In this thesis, we work towards formally proving the correctness of Nipkow's unification algorithm for higher-order lambda patterns \cite{Nipkow-LICS-93} in Isabelle\slash HOL. We use Grobauer's more formal development of the same algorithm \cite{Grobauer} as a basis. We formalize the proofs of most of the auxiliary lemmas used by Grobauer. We discover that for a few of these lemmas the formal proof requires significantly more work and even additional assumptions which are not mentioned explicitly in Grobauer's thesis. Particularly, we find that one of the lemmas requires the lambda terms in question to be simply typed, and we present an untyped lambda term which serves as a counterexample. This extends to the unification algorithm which will not necessarily find the most general unifier if we consider untyped terms. We also discuss how the unification algorithm can be implemented in Isabelle.

Constructing a formal proof in a proof assistant such as Isabelle has the advantage that we can be quite certain that the proof actually works. The proof assistant will check that every step of the proof is justified by logical axioms or by previously proved rules. If we trust the proof assistant and we trust that the formalization expresses what we wanted to express, then we can be certain that the statement we proved is correct. Additionally, using a proof assistant can help with finding mistakes in the argumentation of a pen-and-paper proof. On the other hand, constructing proofs in a proof assistant can be cumbersome because every concept has to be precisely defined, and even though the proof assistant will apply simple proof steps on its own, formal proofs often have to be very elaborate.

We encountered both the advantages and the disadvantages of formalization in this thesis. Translating existing proofs to Isabelle was sometimes very time consuming, but as mentioned above, we were able to find some issues in the pen-and-paper proofs we formalized.

Unfortunately, due to time constraints, it was not possible to formally prove the correctness of the unification algorithm, so this thesis can be seen as a starting point towards a formal proof of correctness or as an extension of the existing implementation of the lambda calculus in the Isabelle\slash HOL distribution.

Chapter \ref{ch:pre} introduces the concepts and notation necessary to understand the rest of the thesis. Chapter \ref{ch:types} explains why types are necessary to obtain most general unifiers. Chapter \ref{ch:lam} describes the formalization of the lambda calculus, based on the existing implementation in the Isabelle\slash HOL distribution and on Grobauer's thesis \cite{Grobauer}. Chapter \ref{ch:unif} presents the unification algorithm and its implementation in Isabelle and the first steps towards proving termination. Finally, Chapter \ref{ch:con} summarizes the thesis and provides an outlook.

\paragraph{Related Work}
The unification algorithm discussed in this thesis was proposed by Nipkow \cite{Nipkow-LICS-93}. Grobauer \cite{Grobauer} presented a more formal development of the same algorithm which we use as a basis. Huet \cite{Huet_1973} showed that general higher-order unification is undecidable and presented a semi-decision algorithm \cite{HUET197527}. Miller \cite{10.1007/BFb0038698} discovered higher-order patterns and showed that unification of higher-order patterns is a decidable subset of general higher-order unification.